% status / progress bar
\section*{Current Status}
\addcontentsline{toc}{section}{Current Status}

\begin{tcolorbox}[title=Status as of \today, colback=green!5, colframe=green!40!black]
\textbf{Active layer:} Layer 0 (\autoref{ch:layer0}), verification mode against the exact eigenmode solution. Code edits to \texttt{heat\_solver.cu} are fully specified; awaiting the seven verification-run CSV rows (three-point spatial sweep, four-point temporal sweep) before Layer 1 opens.

\textbf{Not yet started:} Layers 1--10 (\autoref{ch:layer1}--\autoref{ch:layer10}). Each of those chapters is a scoped placeholder -- the plan was agreed in advance, but no code or results exist for them yet. An empty chapter body means exactly that, not that the layer was skipped or forgotten.

\textbf{Open, unresolved methodological questions:} see \autoref{app:process} for the full tracker -- most importantly, what ``equivalent fault'' means across the classical and quantum substrates (\autoref{app:fault}), and whether injection timing should be tied to real NMDB cosmic-ray flux data.

\textbf{For visiting readers (Sergio, Ricard):} this document tracks the project chapter-by-chapter as each layer is actually built, rather than being written as a finished paper. Please read status markers literally.
\end{tcolorbox}